\problemname{Stökiga känguruungar}

A \href{https://en.wikipedia.org/wiki/Kangaroo_word}{kangaroo word} is a word that carries a synonym of itself (a ``joey''), in such a way that all of the synonym's letters appear in the word, in the same order.
For example, \texttt{pastej} is a kangaroo word, since it carries the synonym \texttt{paj} (\texttt{\textbf{pa}ste\textbf{j}}).
Also \texttt{aste} and \texttt{atj} would count as joeys if we pretend they were words, but not \texttt{paaj} or \texttt{etsa}.
Formally, the joey must be a \emph{subsequence} of the word.

Furthermore, we say that a joey is \emph{rowdy} if it fits in the word in two different ways.
\texttt{paj} is not a rowdy joey, but if the original word had been \texttt{paastej} it would have been --
then it could have been hidden as either \texttt{\textbf{pa}aste\textbf{j}} or \texttt{\textbf{p}a\textbf{a}ste\textbf{j}}.

Given a (made-up) word $S$, and a list of (made-up) synonyms, how many of the synonyms are rowdy joeys of $S$?

\section*{Input}
\begin{itemize}
  \item
    The first line contains a non-empty string consisting of letters \texttt{a-z}, the word $S$ that we are wondering about.

  \item
    The second line contains the integer $N$ ($1 \le N \le 100\,000$): the number of synonyms of the word.

  \item
    The following $N$ lines contain the synonyms, each a non-empty string consisting of letters \texttt{a-z}.
\end{itemize}

No synonym will appear twice, or be equal to $S$.

Let $M$ denote the number of letters in $S$, and $K$ the sum of the number of letters in the synonyms.
Then $M \le 100\,000$, $K \le 500\,000$.

\section*{Output}
Print a single integer -- the number of words that are rowdy joeys of $S$.

\section*{Scoring}
Your solution will be tested on a set of test groups.
To earn points for a group, you must pass all test cases in that group.

\noindent
\begin{tabular}{| l | l | l |}
  \hline
  Group & Points & Constraints \\ \hline
  $1$   & $20$       & $N, M, K \le 100$. \\ \hline
  $2$   & $20$       & $N, M \le 1000$, and all joeys are rowdy joeys. \\ \hline
  $3$   & $20$       & $N, M \le 1000$. \\ \hline
  $4$   & $20$       & All joeys are rowdy joeys. \\ \hline
  $5$   & $20$       & No additional constraints. \\ \hline
\end{tabular}

\section*{Explanation of Samples}
In sample 1, the first three words are joeys of $S$, and moreover rowdy. The test case could therefore appear in test group 2 or 4.

In sample 2, the first four words are joeys, of which the first two are also rowdy joeys. This test case could not appear in test group 2 or 4.
